% SPDX-License-Identifier: Apache-2.0
\section{A Counted Step in a Process of Several Waits}
\label{sec:x-serial}

Section~\ref{sec:x-seq} numbered a loop's waits into states. A step
that repeats a fixed number of times, a bit a cycle for eight cycles,
is a \code{for} whose body waits, and the lowering unrolls it: a
state per turn, the loop's variable a number in each. So
\code{for i in 0..8} around one wait is eight states, and
\code{bit(i)} in them is \code{bit(0)} to \code{bit(7)}. A
\code{for} inside a \code{for} unrolls the same way, inner first.

The bounds are written out, \code{0..8} or \code{0..=7}, because the
macro counts the states and cannot see a const parameter. A
\code{for} that does not wait keeps the unrolling of
Section~\ref{sec:x-loop}, which happens when \code{lowered} runs and
so takes a bound the macro cannot see; the two are told apart by
whether the body waits. A wait inside a \code{for} was dropped
without a word before this, one clocked block doing every turn at
once (issue 587), and is now the states it means.

The serialiser below takes a byte and sends its bits, low bit first,
one a cycle. Nine states: the wait for the byte and the eight turns.
Its netlist shows what the unrolling did: the bit channel's data is
the bit the state selects, and its valid is the eight turns. The
build checks the netlist against the run under nvc and Verilator.

\begin{figure}[htbp]
\centering
\input{serial_timing.tex}
\caption{The serialiser: a byte is taken, and its eight bits follow
one a cycle, low bit first; the second byte waits until the first is
out.}
\label{fig:x-serial}
\end{figure}

\lstinputlisting[caption={\code{ex\_serial.rs}. Run with \code{bazel run //lib/examples:ex\_serial}.},label={lst:x-serial}]{lib/examples/ex_serial.rs}

\lstinputlisting[style=ascii,caption={What \code{ex\_serial} printed when the build ran it: the bits as they come, then the Verilog, one state per turn of the \code{for}.},label={lst:x-serial-out}]{out_serial.txt}
